% Preamble
\documentclass[../main.tex]{subfiles}

% Document
\begin{document}
\chapter{State of the art}\label{ch:stateofart}

In this section an explanation will be given of what both the multi-armed bandit
problem and the contextual multi-armed bandit problem consist of, and the main
algorithms that have been used classically in the literature to address these problems
will be explained in detail.  

The multi-armed bandit problem arises when one is presented with n different options
(or actions) repeatedly and intends to maximize the reward obtained from the chosen
options over a stipulated period of time.  

These rewards usually consist of numerical values coming from a stationary probability
distribution that differs according to the action chosen.  

The example most commonly used to explain this problem is that of a set of slot
machines. Each slot machine translates into a different arm of the multi-armed bandit
and the objective of the problem will be to decide which machine to choose to maximize
the amount of money obtained from the machines' prizes. To maximize these prizes, one
will try to choose those slot machines that have a higher probability of giving a
prize.  

Each of the actions of a multi-armed bandit has an expected average reward, which, if
they were known for each one of the arms, would turn this problem into a trivial
problem, since one would always choose the action with the best average reward.
However, these expected average rewards are not known in most cases, or at least not
exactly. For this reason, in order to estimate what the average reward of each of the
actions is, it will be necessary to try choosing each one of them and to evaluate the
rewards obtained so that an estimate can be made of the average of each possible
option.  

Obtaining a good balance between trying the different possible actions to improve the
estimates of the average rewards of each of these actions and choosing the action with
the best average reward according to the available estimates to maximize the reward
obtained is something fundamental in this type of problems, and is what is known as
the exploration versus exploitation dilemma or conflict.  

This phenomenon of exploration vs. exploitation is a concept widely studied in the
field of reinforcement learning, and there is no unique solution that allows solving
this problem for any multi-armed bandit.  

For this reason, there exists a large number of algorithms that are used nowadays to
solve this class of problems, each of them balancing in a different way this
equilibrium between exploration and exploitation.  

Generally speaking, algorithms with a greater degree of exploitation will obtain
better accumulated rewards in short periods of time, while algorithms with a greater
degree of exploration will obtain better accumulated rewards in longer periods of
time, since the latter are more likely to find actions with better average rewards
than the actions that at first were believed to have the best average rewards among
all possible actions.  

The most commonly used classical algorithms for solving multi-armed bandit problems
are the $\epsilon$-greedy algorithm, Upper Confidence Bound (UCB), and Thompson
Sampling (TS). Below, the mathematics behind each of these algorithms will be
explained. For more detailed information, \textcite{sutton2018} can be consulted for
the first two algorithms, and \textcite{russo2017} for Thompson Sampling.  

\textbf{$\epsilon$-greedy algorithm}: This algorithm or group of algorithms consists
of choosing the action that is expected to give the best reward with probability 
$\left(1 - \epsilon\right)$. From a formal point of view one has that:  

Given an action $a$, the estimate of the value of said action at a time step $t$ will
be  

\begin{equation*}
    Q_{t}\left(a\right) = \dfrac{r_{1} + r_{2} + ... + r_{n_{t}\left(a\right)}}{n_{t}\left(a\right)}
\end{equation*}

Where $r_{1}, r_{2}, ..., r_{n_{t}\left(a\right)}$ are the rewards of the 
$n_{t}\left(a\right)$ previous time steps in which action a was chosen. This method 
to estimate the value of the actions is known as the sample-average method.  

When $n_{t}\left(a\right) = 0$, a default value is assigned to $Q_{t}\left(a\right)$,
and when $n_{t}\left(a\right) \rightarrow \inf$, by the law of large numbers one has
that $Q_{t}\left(a\right) = q\left(a\right)$, where $q\left(a\right)$ is the true
value of action $a$.  

A \textit{greedy} algorithm will choose at time $t$ the action:  

\begin{equation*}
    a_{t} = argmax_{a} Q_{t}\left(a\right)
\end{equation*}

Where $a_{t}$ is the action chosen at time $t$. In case there are several actions that
maximize $Q_{t}$, the chosen action will be a random one among all that maximize it.  

The $\epsilon$-greedy algorithm will choose that action with probability 
$\left(1 - \epsilon\right)$, with $\epsilon$ being a fixed value belonging to the
interval $\left[0, 1\right]$, and a different action chosen completely at random with
probability$\epsilon$.  

When $\epsilon = 0$, we will be in the case of a completely \textit{greedy} algorithm,
and when $\epsilon = 1$, we will be in the case of an algorithm that chooses actions
completely at random.  

\textbf{Upper Confidence Bound (UCB) algorithm}: Contrary to the case of 
$\epsilon$-greedy methods, the exploration of the actions that a priori are not
optimal is carried out following a strategy different from choosing a random action
among all possible ones. In the case of the UCB algorithm, the choice of the action to
explore is made depending on the potential that each of the actions has to end up
maximizing the reward, taking into account both how close the estimate of the mean of
each action is to the action whose average reward is the maximum, and the uncertainty
of such estimates.  

The decision of which action to choose is made in this case following the formula:  

\begin{equation*}
    a_{t} = argmax_{a}\left[Q_{t}\left(a\right) + \alpha \sqrt{\dfrac{ln t}{n_{t}\left(a\right)}}\right]
\end{equation*}

Where $ln t$ denotes the natural logarithm of $t$ and $\alpha > 0$ is an adjustable 
parameter that controls the degree of exploration of the algorithm. If 
$n_{t}\left(a\right) = 0$, $a$ is considered to be maximizing the action, thus avoiding 
divisions by 0.  

The term that is found inside the square root serves as a measure of the uncertainty 
of the estimation of the reward value of action $a$.  

\textbf{Thompson Sampling (TS) algorithm}: This algorithm has a more probabilistic 
component than the previous algorithms, which are based on estimating the action to 
choose empirically and directly from the results of the actions chosen in previous 
time steps.  

In this case, it is assumed that the rewards of the different actions to choose from 
follow probability distributions defined a priori and that are updated each time an 
action is chosen and a reward is obtained.  

More formally: given an a priori probability distribution $p_{0}\left(a\right)$, for 
each action $a$, at each step of the Thompson Sampling algorithm a 
$\hat{Q}_{t}\left(a\right) ~ p_{t}\left(a\right)$ is randomly sampled, where 
$p_{t}\left(a\right)$ is the a posteriori probability distribution at step $t$ of 
action $a$. For that time $t$, the chosen action will be:  

\begin{equation*}
    a_{t} = argmax_{a}\hat{Q}_{t}\left(a\right)
\end{equation*}

The a posteriori probability distribution is obtained from $p_{0}\left(a\right)$, 
applying Bayes’ theorem in the following way:  

\begin{equation*}
    p\left(a | r_{1}, ..., r_{n_{t}\left(a\right)}\right) = \dfrac{p\left(r_{1}, ..., r_{n_{t}\left(a\right)} | a\right) \cdot p_{0}\left(a\right)}{p\left(r_{1}, ..., r_{n_{t}\left(a\right)}\right)}
\end{equation*}

Where $a$ is the action for which the probability distribution is to be estimated, 
$r_{1}, ..., r_{n_{t}\left(a\right)}$ are the rewards of the $n_{t}\left(a\right)$ 
previous time steps in which said action was chosen, and 
$p\left(a | r_{1}, ..., r_{n_{t}\left(a\right)}\right) = p_{t}\left(a\right)$.  

This update of the a posteriori distribution can be carried out at each time step 
by means of the recursive formula:  

\begin{equation*}
    p_{t+1}\left(a\right) \propto p\left(r_{t} | a\right) \cdot p_{t}\left(a\right)
\end{equation*}

A widely known and studied extension of the multi-armed bandit problem is the 
so-called contextual multi-armed bandit. In this problem, in addition to the 
information available in a classical multi-armed bandit problem, there is also a context, 
which is known before making any decision and provides relevant information about the 
environment or the user, potentially helping to make a better decision regarding which 
action to choose.  

This context, which will be denoted as $x_{t} \in X$, is usually given as a feature vector. 
The different available contexts are used to fit models (a different model for each 
available action) to predict the reward that will be obtained in each case.  

Classically, the models used for this purpose were linear models, although nowadays more 
complex nonlinear models based on neural networks or decision trees are also used.  

In this document, the Linear Upper Confidence Bound (LinUCB) and Linear Thompson Sampling 
(LinTS) algorithms will be covered, which are extensions of the previously presented 
algorithms to contextual multi-armed bandits through the use of linear algorithms to model 
the context. For more information about the first algorithm, see \textcite{zhou2015}, and 
for the second, \textcite{russo2017}.  

\textbf{Linear Upper Confidence Bound (LinUCB) algorithm}: It consists of an extension of 
the UCB algorithm through the use of linear models to represent the context. More 
rigorously, this algorithm can be stated as follows:

Given a context $x_{t, a} \in \mathbb{R}^{d}$ for each action $a$ and each time $t$, the 
expected reward is assumed to be a linear function of the context in the following form:  

\begin{equation*}
    Q_{t}\left(a\right) = x_{t, a}^{\top}\theta^{*}
\end{equation*}

Where $\theta^{*} \in \mathbb{R}^{d}$ is an unknown parameter vector which allows 
minimization of the loss function  

\begin{equation*}
    R\left(T\right) = \sum_{t = 1}^{T} \left(x_{t, a_{t}^{*}}^{\top}\theta^{*} - x_{t, a_{t}}^{\top}\theta^{*}\right)
\end{equation*}

Where $a_{t}^{*} =argmax_{a \in A} x_{t, a}^{\top}\theta^{*}$.  

This loss function is known as cumulative regret.  

At each step of the algorithm, there is an estimate of $\theta^{*}$ obtained from 
the linear model. This estimate will be denoted as $\hat{\theta}_{t}$, and the 
action chosen at each step will be:  

\begin{equation*}
    a_{t} = argmax_{a} x_{t, a}^{\top}\hat{\theta}_{t-1} + \alpha \cdot \sqrt{x_{t, a}^{\top} A_{t-1}^{-1}x_{t, a}}
\end{equation*}

Where $\alpha > 0$ is the same adjustable parameter that controls the degree of 
exploration of the UCB algorithm, and the matrix $A_{t-1}$ comes from the linear 
model, which consists of a ridge (regularized) regression of the following form:  

\begin{equation*}
    \hat{\theta}_{t} = A_{t}^{-1} b_{t}
\end{equation*}

Where  

\begin{equation*}
    A_{t} = \lambda I + \sum_{s=1}^{t} x_{s, a_{s}} x_{s, a_{s}}^{\top} \in \mathbb{R}^{d \times d}
\end{equation*}

\begin{equation*}
    b_{t} = \sum_{s=1}^{t} r_{s} x_{s, a_{s}} \in \mathbb{R^{d}}
\end{equation*}

\begin{equation*}
    A_{0} = \lambda I
\end{equation*}

with $\lambda$ being the trainable parameter of the linear regression, $I$ the 
identity matrix, $r_{s}$ the reward at time $s$, and $x_{s, a_{s}}$ the context 
for action $a_{s}$ at time $s$.  

\textbf{Linear Thompson Sampling (LinTS) algorithm}: It is an extension of the 
Thompson Sampling algorithm through the use of linear models to account for 
context.  

As in the previous algorithm, the expected reward is given by the equation:  

\begin{equation*}
    Q_{t}\left(a\right) = x_{t, a}^{\top}\theta^{*}
\end{equation*}

However, in this case, $\hat{\theta}_{t}$, the estimate of the parameter vector 
$\theta^{*}$, consists of a random sample from a posterior distribution over 
parameter vectors, $p_{t}\left(\theta\right)$, which is computed from a prior 
distribution $p_{0}\left(\theta\right)$, applying Bayes’ theorem in a manner 
analogous to the Thompson Sampling algorithm, provided that this prior 
distribution is conjugate to the posterior, i.e., they belong to the same family. 
If this is not the case, numerical methods such as Markov Chain Monte Carlo (MCMC) 
must be used to compute the posterior distribution.  

Therefore, taking $\hat{\theta}_{t} \sim p_{t}\left(\theta\right)$, the chosen 
action is:  

\begin{equation*}
    a_{t} = argmax_{a}x_{t, a}^{\top}\hat{\theta}_{t}
\end{equation*}

Finally, it is worth mentioning that extending these linear models to incorporate 
context into nonlinear models consists of replacing the linear expression 
$Q_{t}\left(a\right)=x_{t, a}^{\top} \theta^{*}$ with a nonlinear expression 
$Q_{t}\left(a\right)=f^{*}\left(x_{t , a}\right)$, where $f^{*}$ is a nonlinear 
function.  

However, this transition from a linear function to a nonlinear one entails a 
series of theoretical complications that will not be addressed in this work.  
\end{document}